\documentclass[11pt, a4paper]{awesome-cv}
\geometry{left=1.8cm, top=1.4cm, right=1.8cm, bottom=1.8cm, footskip=.5cm}
\fontdir[fonts/]

\colorlet{awesome}{awesome-skyblue}
\renewcommand{\acvHeaderSocialSep}{\quad\textbar\quad}

\name{Carlo}{Nicolini, Ph.D.}
\position{Senior / Principal AI Research Scientist \textbar{} Statistical Physics \& Frontier AI Interpretability}
\address{Milan / Rovereto, Italy \textbar{} Open to Global Fully Remote Roles}


\homepage{carlonicolini.github.io}
\googlescholar{jnpIfCwAAAAJ}{Google Scholar (800+ citations, h-index 11)}
\github{carlonicolini}
\linkedin{carlo-nicolini}

\makecvfooter{\today}{Carlo Nicolini, Ph.D.~~~·~~~Curriculum Vitae}{\thepage}

\begin{document}
\makecvheader

\cvsection{Executive Summary}

Senior AI Research Scientist and theoretical physicist with 15+ years of experience bridging statistical mechanics, information theory, and frontier machine learning. Currently leading research at Ipazia SpA on Large Language Model interpretability, Vector Symbolic Architectures (TMLR 2026), long-horizon planning benchmarks (SokoBench, TMLR 2026), and inference-time reasoning scaffolds. Maintainer of \textbf{skfolio} (portfolio optimization) and Board Member at \textbf{skfolio Labs}; Fellow of the \textbf{Thinking About Thinking Society}. Published 20+ peer-reviewed papers across TMLR, COLM, EPJ Data Science, Physical Review E, and NeuroImage (800+ citations, h-index 11). Track record of transforming abstract mathematical principles into auditable, high-performance research infrastructure for frontier AI labs.



\cvsection{Experience}



\begin{cventries}

\cventry
{Senior AI Research Scientist} {Ipazia SpA} {Milan, Italy (Remote)} {2022 - Present} { \begin{cvitems}
\item{Lead fundamental and applied AI research on modern Large Language Models (LLMs), mechanistic interpretability, and compound AI systems.}
\item{Pioneered the \textbf{Hyperdimensional Probe} (TMLR 2026), connecting Vector Symbolic Architectures (VSA) with neural activations to extract human-interpretable concept spaces and latent mechanics across foundation models.}
\item{Developed \textbf{SokoBench} (TMLR 2026), an automated evaluation suite testing long-horizon planning, tree search, and reasoning frontiers in state-of-the-art LLMs against verified combinatorial solvers.}
\item{Engineered state-of-the-art information extraction and auditable RAG architectures for enterprise sustainability disclosure, published in \textit{EPJ Data Science} (110+ citations).}
\item{Designed mathematically grounded inference-time scaffolds (branching, verification, and probabilistic language programming), turning theoretical physics principles into reliable, production-grade AI systems.}
\end{cvitems}
}

\cventry
{Co-Author \& Core Maintainer \textbar{} Board Member} {skfolio \& skfolio Labs} {Open Source \textbar{} skfoliolabs.com} {2022 - Present} { \begin{cvitems}
\item{Maintain and drive architectural design for \textbf{skfolio} (\url{https://skfolio.org}), the leading scikit-learn compatible open-source library for mathematical portfolio optimization and risk modeling.}
\item{Engineered high-performance modules for Online Convex Optimization (regret minimization), compiled solver acceleration (\textit{skfolio-accelerate}), and autonomous LLM-driven strategy discovery (\textit{autoresearch-skfolio}).}
\item{Member of the Board at \textbf{skfolio Labs}, steering scientific strategy, open-source governance, and enterprise partnerships.}
\item{Designed and published novel optimization methodologies including Hopfield Networks for discrete asset allocation presented at \textit{ACM ICAIF 2024}.}
\end{cvitems}
}

\cventry
{Senior Quantitative Data Scientist \& Data Engineer} {Prometeia} {Milan, Italy} {2020 - 2022} { \begin{cvitems}
\item{Spearheaded quantitative modeling and large-scale distributed analytics (Apache Spark, Python) for systemic financial risk and non-performing loan (NPL) forecasting.}
\item{Designed high-throughput, graph-based anti-money laundering (AML) detection pipelines combining complex network algorithms with interpretable supervised ML.}
\item{Bridge technical research and executive delivery under strict regulatory, risk, and institutional audit constraints.}
\end{cvitems}
}

\cventry
{Postdoctoral Research Fellow (Statistical Physics \& Complex Networks)} {Istituto Italiano di Tecnologia (IIT)} {Rovereto, Italy} {2017 - 2020} { \begin{cvitems}
\item{Formulated quantum-inspired and statistical mechanics frameworks (spectral von Neumann entropy, density matrices) for multi-scale network analysis and brain functional connectivity.}
\item{Authored high-impact publications in \textit{Physical Review E}, \textit{NeuroImage}, and \textit{Scientific Reports}, accumulating over 800 citations and establishing novel resolutions to classical graph community detection limits.}
\item{Authored and distributed open-source scientific software packages (\textbf{Networkqit}, \textbf{PACO}, \textbf{xyz}) in C++ and Python for high-performance network thermodynamics and information-theoretic estimation.}
\end{cvitems}
}

\cventry
{Scientific Programmer \& Research Assistant} {Istituto Italiano di Tecnologia \& University of Trento} {Rovereto \& Trento, Italy} {2009 - 2017} { \begin{cvitems}
\item{Designed real-time 3D stereoscopic tracking and computational psychophysics platforms (C++, OpenGL, CMake).}
\item{Conducted research on machine learning optimization, kernel methods (SVMs), and high-dimensional human behavioral tracking.}
\end{cvitems}
}

\end{cventries} 

\cvsection{Technical \& Research Leadership}



\begin{cvskills}

\cvskill
{AI Interpretability \& Reasoning} {Mechanistic interpretability, Vector Symbolic Architectures (VSA), linear / hyperdimensional probing, latent representation geometry, inference-time reasoning scaffolds, long-horizon planning evaluation (SokoBench), compound AI systems, probabilistic language programming, auditable RAG architectures.} 

\cvskill
{Statistical Physics \& Math} {Statistical field theory of complex systems, non-equilibrium dynamics, spectral von Neumann entropy, density matrices, graph Laplacians, Wilson's random forest sampling, community detection / resolution limits, online convex optimization, regret minimization, information theory.} 

\cvskill
{Scientific Software \& Systems} {Python (NumPy, SciPy, PyTorch, scikit-learn, JAX, Hugging Face, Numba), C++ (C++17/20, Eigen, Boost, CMake), Rust bindings, distributed compute (Ray, Spark), Git, Linux, CI/CD, production-grade API architecture \& open-source maintenance.} 

\cvskill
{Quantitative Modeling} {Modern portfolio theory, risk budgeting, convex optimization (CVXPY), walk-forward cross-validation, financial time series, Hopfield networks for combinatorial allocation, anti-money laundering (AML) network graphs.} 

\end{cvskills} 

\cvsection{Selected Publications}



\begin{cventries}

\cventry
{Co-Author} {Sokobench: Evaluating Long-Horizon Planning and Reasoning in Large Language Models} {Transactions on Machine Learning Research (TMLR)} {2026} {Sebastiano Monti, \textbf{Carlo Nicolini}, Gianni Pellegrini, Jacopo Staiano, Bruno Lepri}

\cventry
{Co-Author} {Hyperdimensional Probe: Decoding LLM Representations via Vector Symbolic Architectures} {Transactions on Machine Learning Research (TMLR)} {2026} {Marco Bronzini, \textbf{Carlo Nicolini}, Bruno Lepri, Jacopo Staiano, Andrea Passerini}

\cventry
{Co-Author} {Unveiling LLMs: The Evolution of Latent Representations in a Dynamic Knowledge Graph} {Conference on Language Modeling (COLM)} {2025} {Marco Bronzini, \textbf{Carlo Nicolini}, Bruno Lepri, Jacopo Staiano, Andrea Passerini}

\cventry
{Lead Author} {Hopfield Networks for Asset Allocation} {ACM International Conference on AI in Finance (ICAIF)} {2024} {\textbf{Carlo Nicolini}, Monisha Gopalan, Jacopo Staiano, Bruno Lepri}

\cventry
{Lead Author} {Spectral Entropy of Random Spanning Forests} {arXiv preprint (Thermodynamic Sampling via Wilson Algorithm)} {2025} {\textbf{Carlo Nicolini}}

\cventry
{Co-Author} {Glitter or Gold? Deriving Structured Insights from Sustainability Reports via Large Language Models} {EPJ Data Science (110+ citations)} {2024} {Marco Bronzini, \textbf{Carlo Nicolini}, Bruno Lepri, Jacopo Staiano, Andrea Passerini}

\cventry
{Co-Author} {Statistical Physics of Complex Information Dynamics} {Physical Review E (70+ citations)} {2020} {Arsham Ghavasieh, \textbf{Carlo Nicolini}, Manlio De Domenico}

\cventry
{Lead Author} {Thermodynamics of Network Model Fitting with Spectral Entropies} {Physical Review E} {2018} {\textbf{Carlo Nicolini}, Vladimir Vlasov, Angelo Bifone}

\cventry
{Lead Author} {Scale-Resolved Analysis of Brain Functional Connectivity Networks with Spectral Entropy} {NeuroImage (50+ citations)} {2020} {\textbf{Carlo Nicolini}, Giulia Forcellini, Ludovico Minati, Angelo Bifone}

\cventry
{Lead Author} {Modular Structure of Brain Functional Networks: Breaking the Resolution Limit by Surprise} {Scientific Reports (100+ citations)} {2016} {\textbf{Carlo Nicolini}, Angelo Bifone}

\cventry
{Co-Author} {Graph Analysis and Modularity of Brain Functional Connectivity: Searching for the Optimal Threshold} {Frontiers in Neuroscience (170+ citations)} {2017} {Cécile Bordier, \textbf{Carlo Nicolini}, Angelo Bifone}

\end{cventries} 

\cvsection{Affiliations \& Leadership}



\begin{cvskills}

\cvskill
{Thinking About Thinking} {Fellow of the \textbf{Thinking About Thinking Society}, an international interdisciplinary research group (Oxford) focused on the intersection of mathematics, theoretical neuroscience, cognition, and AI reasoning foundations.}

\cvskill
{skfolio Labs} {Board Member and Scientific Advisor at \textbf{skfolio Labs} (\url{https://skfoliolabs.com}), guiding research strategy, open-source governance, and high-performance algorithmic development for quantitative investment.}

\cvskill
{Scientific Open Source} {Core maintainer of \textbf{skfolio} (portfolio optimization); creator of \textbf{Networkqit} (quantum information theory for networks), \textbf{xyz} (information-theoretic estimators), \textbf{loop-erased-thermodynamics}, and \textbf{PACO}.}

\cvskill
{Peer Review \& Service} {Reviewer for top machine learning and physical sciences venues (TMLR, Physical Review, NeuroImage, Scientific Reports).}

\end{cvskills}
 

\cvsection{Education}



\begin{cventries}

\cventry
{Ph.D. in Applied Physics (Nanoscience \& Complex Systems)} {University of Verona} {Verona, Italy} {2014 - 2017} { \begin{cvitems}
\item {Doctoral dissertation: \textit{The modular structure of brain functional connectivity networks: a graph theoretical approach}. Developed novel information-theoretic and Surprise-based algorithms overcoming fundamental resolution limits in network community detection.}
\item {Conducted at the Center for Neuroscience and Cognitive Systems (CNCS), Istituto Italiano di Tecnologia (IIT).}
\end{cvitems}
}

\cventry
{M.Sc. in Physics and Biomedical Technologies \textbar{} 110/110 cum laude} {University of Trento} {Trento, Italy} {2006 - 2008} { \begin{cvitems}
\item {Thesis on Monte Carlo simulations and stochastic optimization for inverse radiation treatment planning.}
\end{cvitems}
}

\cventry
{B.Sc. in Physics \textbar{} 106/110} {University of Trento} {Trento, Italy} {2003 - 2006} { \begin{cvitems}
\item {Thesis on nonlinear dynamics and feedback systems in cardiovascular pressure regulation.}
\end{cvitems}
}



\end{cventries} 
\end{document}